\documentclass[10pt,a4paper]{amsart}
\usepackage[margin=19mm]{geometry}
\usepackage{amsmath,amssymb,mathtools}
\usepackage[scr=rsfs,cal=euler]{mathalfa}
\usepackage{xfrac}
\usepackage[unicode,hidelinks,bookmarks=false]{hyperref}
\newcommand{\ie}{i.e.\ }
\newcommand{\cf}{cf.\ }
\newcommand{\resp}{respectively\ }
\newcommand{\Subsection}{\subsection}
\providecommand{\nobreakdash}{\nobreak\hbox{-}}
\setlength{\emergencystretch}{2em}
\multlinegap0pt
\usepackage{bm}
\usepackage{bbm}
\usepackage{dsfont}
\usepackage{xspace}
\usepackage{cancel}
\usepackage{mathtools}
\usepackage{amsthm}
\makeatletter
\@ifundefined{theorem}{\newtheorem{theorem}{Theorem}}{}
\@ifundefined{lemma}{\newtheorem{lemma}[theorem]{Lemma}}{}
\makeatother
\makeatletter
\newcommand{\Rr}{\mathbf{R}}
\newcommand{\Zz}{\mathbf{Z}}
\expandafter\ifx\csname epigraph\endcsname\relax
\newenvironment{epigraph}
{\hfill\begin{minipage}{0.6\linewidth}\raggedleft\footnotesize}{\end{minipage}\bigskip\bigskip}
\fi
\newcommand{\expect}{\mathbf{E}}
\renewcommand{\rho}{\varrho}
\DeclareMathOperator{\wass}{\mathcal{W}}
\makeatother
\title[Wasserstein metrics and quantitative equidistribution of exp]{Wasserstein metrics and quantitative equidistribution of exponential sums over finite fields}
\author{Emmanuel Kowalski, Th\'eo Untrau}
\date{Source: arXiv:2505.22059v2 (CC BY 4.0)}
\begin{document}
\maketitle

\noindent\emph{Source and licence.} Wasserstein metrics and quantitative equidistribution of exponential sums over finite fields. Version arXiv:2505.22059v2 (CC BY 4.0), distributed under the Creative Commons Attribution 4.0 International licence, as recorded in the arXiv licence metadata for this exact version and shown on the abstract page https://arxiv.org/abs/2505.22059v2. Full rights notice: licenses/arxiv-2505.22059-CC-BY-4.0.txt.\par\smallskip

\noindent\textit{Editorial scope.} Counted unit: the Lemma whose source label is \texttt{lem: choice of H} in the article (source0.tex lines 2813--2818, source label \texttt{lem: choice of H}), delivered together with the article's own complete original proof of it (source0.tex lines 2820--2870, one \texttt{proof} environment of 51 lines). No mathematics is written, repaired or re-derived: statement and proof are the article's own text line for line. Required context retained uncounted: lem: regularization cost (lines 2772--2781). Every definition, display and label of those lines is retained unchanged. Omitted from this package and not redistributed: 1--2771, 2782--2812, 2819--2819, 2871--2934. Nothing inside a retained excerpt is omitted or altered: every retained body is delivered exactly as the article's lines are written, so no omission receipt of any kind accompanies this package. Citations: the retained excerpts carry no citation command at all, so no bibliography entry of the article is transcribed and no reference is redistributed. Reference adaptation: none was needed and none was made; every reference inside the delivered unit resolves inside it, so no unresolved reference or citation is produced. Printed slips kept literal: no printed word, symbol, hypothesis or number of the counted statement or of its proof was altered. Layout and class: the article's own class is \texttt{amsart} with packages including amsxtra, amssymb, amsthm, amsmath, amscd, url, amsrefs, enumitem, inputenc, eucal, fullpage, mathrsfs, csquotes, hyperref; the wrapper is the lane's pinned \texttt{amsart} editorial wrapper at 10pt on A4, which restates the theorem-like environments the retained text needs and adds the article's own macro definitions that the retained text needs (\texttt{\textbackslash Rr}, \texttt{\textbackslash Zz}, \texttt{\textbackslash expect}, \texttt{\textbackslash rho}, \texttt{\textbackslash wass}), with extra wrapper packages (bm, bbm, dsfont, xspace, cancel, mathtools). The delivered numbering is the unit's own. Assets: no retained span contains a figure environment, a graphics-inclusion command, a PDF-page inclusion, a table or a TikZ picture; no image file is copied into this package, the package declares no asset dependency and no image file is redistributed with it. Licence: version arXiv:2505.22059v2 of this article is distributed under the Creative Commons Attribution 4.0 International licence, as recorded in the arXiv licence metadata for this exact version and shown on the abstract page https://arxiv.org/abs/2505.22059v2. A full rights notice is delivered at licenses/arxiv-2505.22059-CC-BY-4.0.txt. Characters: every retained excerpt is pure ASCII, so the delivered engine, XeLaTeX with its default Latin Modern text font, reports no missing character.
\begin{lemma}\label{lem: regularization cost}
  Let $\mu$ and $\nu$ be two Borel probability measures on
  $(Q^d, \rho_d)$ and let $H = (H_1, \dots, H_d)$ be a random vector
  in $\Rr^d$. For $x \in \Rr$, denote by $M(x)$ the unique element of
  $(x + 2 \pi \Zz)\cap(-\pi, \pi]$. Let $N$ be the random vector
  $(M(H_1), \dots, M(H_d))$ and let $\eta$ denote the law of $N$. Then
  $$
  \wass_1(\mu, \nu) \leqslant 	\wass_1(\mu* \eta, \nu * \eta) + 2 \expect(|H|).
  $$
\end{lemma}

\begin{lemma}\label{lem: choice of H}
  For all integers $T \geqslant 1$, there exists a random variable $H$
  with values in $\Rr^d$ whose characteristic function is supported on
  the cube $[-T,T]^d$ and such that
  $\expect(|H|) \leqslant \frac{2 \sqrt{3}\sqrt{d}}{T}$.
\end{lemma}

\begin{proof}
  As a first attempt (that will require some adjustments) let us use
  what is known on the Fourier transform of the triangle function: it
  is a classic fact that if $X$ is a real random variable whose
  density function is
  $$
  f(x) = \frac{2(1 - \cos(x/2))}{\pi x^2}
  $$
  then its characteristic function is given by
  $$
  \varphi_X(s) = (1 - 2 |s|)^{+}
  $$
  where $(-)^{+}$ denotes the positive part. In particular,
  $\varphi_X$ is supported on $\left[ -\frac12, \frac12 \right]$, so
  that suitable renormalizations of $X$ will allow us to construct
  random variables with support in $[-T, T]$ for all $T$. However, the
  issue is that $\expect(|X|) = + \infty$, which would make the
  inequality of Lemma \ref{lem: regularization cost} useless. This is
  why the following adjustment is needed: rather than working with
  $X$, we will work with the random variable $\xi$ whose
  characteristic function is given by
  \begin{equation} \label{eq: def v*v}
    w(s) = 3 (\varphi_X * \varphi_X) (s) = 3 \int_{\Rr} \varphi_X(s-t)\varphi_X(t) dt.
  \end{equation}
  Here, the factor $3$ is just a normalization factor to ensure that
  $w(0) = 1$, which is a necessary condition in order to be a
  characteristic function. Then thanks to the convolution theorem for
  the Fourier transform, $w$ is the characteristic function
  $\varphi_\xi$ of a random variable $\xi$ with density
  $$
  g(x) =6 \pi f(x)^2 = \frac{24 (1 - \cos(x/2))^2}{\pi x^4}\cdot
  $$
  Since $w = 3(\varphi_X * \varphi_X)$, it is supported on $[-1,1]$,
  and this time $\expect(|\xi|) < + \infty$. Even better,
  $\expect(\xi^2)$ is also finite, and can be explicitly computed!
  Indeed, since $w$ is the characteristic function of $\xi$, the
  second moment of $\xi$ is equal to $-w''(0)$, and this can be
  calculated from \eqref{eq: def v*v}, yielding $\expect(\xi^2) =
  12$. To conclude, it suffices to define $H$ as
  $\frac{1}{T}(\xi_1, \dots, \xi_d)$ for independent random variables
  $\xi_i$ having the same distribution as $\xi$. Indeed,
  $\expect(|H|^2) = \frac{12d}{T^2}$ and
  $(\expect|H|)^2 \leqslant \expect(|H|^2)$ so
  $\expect(|H|) \leqslant \frac{2 \sqrt{3}\sqrt{d}}{T}$. Moreover, for all
  $s =(s_1, \dots, s_d) \in \Rr^d$,
  $$
  \varphi_H(s) = \prod_{j = 1}^{d} w(s_j / T)
  $$
  so the choice of a $w$ supported on $[-1, 1]$ implies that
  $\varphi_H$ is supported on $[-T,T]^d$.
\end{proof}

\end{document}
